
\chapterspace{-4}
\chapter[Propiedad de ser cohomológicamente propios los haces...]{Propiedad de ser cohomológicamente propios\\ los haces de conjuntos y\\ los haces de grupos no conmutativos}
\label{XIII}
\marginpar{344}

\begin{center}
{por la Sra. M.\ \textsc{Raynaud}\footnote{A partir de notas inéditas de
A\ptbl\ Grothendieck.}}
\end{center}
\vspace*{1cm}

La presente exposición se propone utilizar la cohomología étale para
generalizar ciertos resultados de~\Ref{IX} y~\Ref{X}. Muestra asimismo cómo
se pueden extender a los haces de grupos no necesariamente conmutativos los
resultados de~SGA~5~II que siguen teniendo sentido para tales haces. Se dan
por conocidas las nociones de cohomología étale expuestas en~SGA~4.

El resultado principal \eqref{XIII.2.4} proporciona un ejemplo importante de
un morfismo no propio $f\colon~U\to~S$ que es \og cohomológicamente propio en
dimensión~$\leq1$\fg, es decir, tal que, para ciertos haces de grupos~$F$
sobre~$U$ (en el sentido de la topología étale), la formación de~$f_*F$ y
$\R^1f_*F$ conmuta con todo cambio de base~$S'\to~S$. En efecto, esta
propiedad la satisface el abierto~$U$ de un esquema~$X$ propio sobre~$S$ que
es el complementario de un divisor~$D$ con cruzamientos normales
relativamente a~$S$, al menos si se exige que~$F$ sea constante finito y de
orden primo con las características residuales de~$S$. Si ya no se supone que
$F$ sea de orden primo con las características residuales de~$S$, se obtiene
un resultado análogo sustituyendo~$\R^1f_*F$ por el subhaz
$\R^1_{\tame}f_*F$ obtenido al limitarse a considerar los torsores bajo~$F$
\og moderadamente ramificados sobre~$X$ relativamente a~$S$\fg. En
particular, esto permite demostrar que el grupo fundamental moderadamente
ramificado de una curva algebraica propia y lisa sobre un cuerpo
separablemente cerrado, privada de un número finito de puntos cerrados, está
topológicamente generado por un número finito de elementos
\eqref{XIII.2.12}.

El
\marginpar{345}
\No \Ref{XIII.4} está dedicado a la sucesión exacta de homotopía y a la
fórmula de Künneth.

Por último, un apéndice ofrece variantes útiles del lema de Abhyankar
demostrado en~\Ref{X.3.6}.

\setcounter{section}{-1}
